\section{Cuts in the original variables}\label{sec:original}

Proposition~\ref{prop:support} produces inequalities in lifted
coordinates $(x,w)$. Using them requires variables $w_j$ with
$w_j=f_j(x)$ in the model. If the solver's representation does not
already contain exactly these variables, adding them changes the model:
new equalities, new variables for branching and propagation, and possibly
a different presolve. Our first campaign (Section~\ref{sec:campaign1})
used such a reformulation and needed a control mode to separate its
effect from that of the cuts. This section shows how to avoid the
auxiliary variables. The nonlinear functions are minimized jointly, and
the rows of the model then eliminate them. The underlying facts are
classical Lagrangian duality; we state them in the form that the
implementation and its certificates need.

\subsection{Aggregating nonlinear rows}\label{sec:aggregation}

Let $v\in\R^N$ be the vector of all model variables, including an
objective variable if there is one, and let $x=Sv$ be a block of $n$ of
them, selected by a $0/1$ matrix $S$. Suppose the model contains rows
\begin{equation}\label{eq:rows}
g(Sv)+Av\le b,\qquad h(Sv)+Cv=e,
\end{equation}
with $g:D\to\R^m$ and $h:D\to\R^p$, and affine remainders $Av$ and $Cv$
that may involve any variables, including block variables. Here $D$ is a
compact set that contains the block part $Sv$ of every point to which the
cuts must apply, for example the block box intersected with the affine
rows that involve only block variables.

\begin{proposition}[Aggregated support cut]\label{prop:aggregate}
Let $a\in\R^n$, $\lambda\in\R^m_{\ge0}$, $\mu\in\R^p$ and
\begin{equation}\label{eq:agg-support}
\beta\ \le\ \varphi(a,\lambda,\mu):=\min_{u\in D}\bigl\{a\T u+\lambda\T g(u)+\mu\T h(u)\bigr\}.
\end{equation}
Then every point that satisfies~\eqref{eq:rows} and $Sv\in D$ satisfies
\begin{equation}\label{eq:agg-cut}
(S\T a-A\T\lambda-C\T\mu)\T v\ \ge\ \beta-\lambda\T b-\mu\T e .
\end{equation}
\end{proposition}

\begin{proof}
At such a point, $\lambda\ge0$ gives $\lambda\T g(Sv)\le\lambda\T(b-Av)$,
and $\mu\T h(Sv)=\mu\T(e-Cv)$. Substitute both into
$a\T Sv+\lambda\T g(Sv)+\mu\T h(Sv)\ge\beta$.
\end{proof}

This is weak Lagrangian duality \citep[Section~5.1]{BoydVandenberghe2004}.
It needs neither convexity nor optimal multipliers. The same inequality is
the Lagrangian cut of a block in the extended block-separable
reformulation of \citet[Section~7.1]{Nowak2005}, after the block's
auxiliary variables are eliminated with the rows~\eqref{eq:rows};
related cuts and range reductions appear in
\citet{TawarmalaniSahinidis2004}, \citet{KaruppiahGrossmann2008} and
\citet{DomesNeumaier2016}. Its strength comes from minimizing the
nonlinear functions jointly over their common domain before they are
eliminated. Nonnegative combinations of linear rows that are already in
the relaxation cannot cut off any point of it.

An objective is handled through its epigraph. For minimization of
$g_0(x)+\ell_0(v)+c_0$ with objective variable $t$, the row
$g_0(x)+\ell_0(v)-t\le-c_0$ is one of the rows~\eqref{eq:rows}; for
maximization the hypograph row $-g_0(x)-\ell_0(v)+t\le c_0$ is used. An
equality may also be written as two opposite inequalities with
nonnegative multipliers, whose difference is the free multiplier.

Validity uses only $Sv\in D$ and the selected rows. Integrality, the
other constraints and the objective play no role, so the cut is valid for
the mixed-integer feasible set and for its closed convex hull. If $D$ is
built from bounds valid only at a node of the search tree, or from an
incumbent cutoff, the cut is valid only there; our implementation uses
only original bounds and rows, so its cuts are global. Integer block
variables can be exploited by taking $D$ discrete, for instance
$D\cap\Z^n$; we do not do this.

\subsection{What the cuts can express}\label{sec:closure}

Let $K=\conv\{(u,g(u),h(u)):u\in D\}$ and let $Q=\{0\}\times\R^m_{\ge0}\times\{0\}$.
The affine map $T(v)=(Sv,\ b-Av,\ e-Cv)$ places a model point in the
lifted space of the block. The next theorem identifies the set described
by all cuts~\eqref{eq:agg-cut}. It is the set version of the classical
fact that the Lagrangian dual of a nonconvex problem equals the dual of
its convexification in the joint space of variables and constraint
values \citep{Falk1969,Geoffrion1974,FeltenmarkKiwiel2000,LemarechalRenaud2001};
the separation argument is the one that \citet[Theorem~3]{ChenLuedtke2022}
use for Lagrangian cuts in two-stage integer programs. We give the proof
for our setting, with nonlinear rows, affine remainders on arbitrary
variables and equality rows, because it is short and needs no constraint
qualification.

\begin{theorem}[Closure of the aggregated cuts]\label{thm:closure}
Let $D\subset\R^n$ be nonempty and compact and let $g$ and $h$ be
continuous on $D$. A point $v$ satisfies every inequality~\eqref{eq:agg-cut}
with $\beta=\varphi(a,\lambda,\mu)$, over all $a\in\R^n$,
$\lambda\in\R^m_{\ge0}$ and $\mu\in\R^p$, if and only if
\begin{equation}\label{eq:closure}
v\in\mathcal C:=T^{-1}(K+Q)=\bigl\{v:\ \exists\,y\ \text{with}\ (Sv,y,e-Cv)\in K\ \text{and}\ y\le b-Av\bigr\}.
\end{equation}
\end{theorem}

\begin{proof}
The set $K$ is compact and convex, and $Q$ is a closed convex cone, so
$U=K+Q$ is closed and convex: if $k_t+q_t\to\bar u$ with $k_t\in K$ and
$q_t\in Q$, a subsequence of $(k_t)$ converges in $K$, and then $(q_t)$
converges in $Q$. A linear functional $(a,\lambda,\mu)$ has finite
infimum on $U$ only if $\lambda\ge0$, and then the infimum is
$\varphi(a,\lambda,\mu)$. Evaluating the valid inequality
$a\T u+\lambda\T y+\mu\T\eta\ge\varphi(a,\lambda,\mu)$ of $U$ at $T(v)$
gives exactly~\eqref{eq:agg-cut}. Hence every point of $\mathcal C$ satisfies all
cuts. Conversely, if $T(v)\notin U$, a linear functional strictly
separates $T(v)$ from $U$; its infimum on $U$ is finite, so
$\lambda\ge0$, and the corresponding cut is violated at $v$.
\end{proof}

If auxiliary variables $w=g(x)$ and $s=h(x)$ were added, the joint-graph
relaxation of the block would be $(x,w,s)\in K$, $w+Av\le b$, $s+Cv=e$,
and its projection onto $v$ is $\mathcal C$. The direct cuts can therefore express
exactly what this lifted relaxation expresses about the original
variables, without adding variables or equalities to the model. They
express no more. Relaxations that keep the rows inside the block, such as
the convex-hull relaxations of block-decomposition methods
\citep[Section~3]{Nowak2005}, \citep{WuMutsNowakHendrix2025}, can be
strictly stronger (Example~\ref{ex:quarter}). There is also a cost: with
the same linearization points, outer approximation of a lifted
formulation can be tighter than outer approximation in the original
space \citep{TawarmalaniSahinidis2005}, and a finite set of direct cuts
may need more members than a lifted description.

For a single row, $\mathcal C$ is the familiar envelope relaxation.

\begin{corollary}[One row]\label{cor:one-row}
If $m=1$ and $p=0$, then $\mathcal C=\{v:\ Sv\in\conv D,\ \breve g(Sv)+Av\le b\}$,
where $\breve g(x)=\min\{y:(x,y)\in K\}$ is the convex envelope of $g$ over
$D$. If $m=0$ and $p=1$, then
$\mathcal C=\{v:\ Sv\in\conv D,\ \breve h(Sv)\le e-Cv\le\hat h(Sv)\}$, where
$\hat h(x)=\max\{y:(x,y)\in K\}$ is the concave envelope.
\end{corollary}

\begin{proof}
The projection of $K$ onto $x$ is $\conv D$, and for each $x\in\conv D$
the fiber $\{y:(x,y)\in K\}$ is the interval $[\breve g(x),\hat g(x)]$. A
suitable $y$ in~\eqref{eq:closure} exists if and only if
$\breve g(Sv)\le b-Av$. The equality case is the same with $y=e-Cv$.
\end{proof}

Joint treatment can only help when rows share block variables or when
$D$ couples them.

\begin{proposition}[No gain without coupling]\label{prop:product}
Suppose $x=(x^1,x^2)$, $D=D^1\times D^2$ with compact factors, and the
rows split into two groups whose nonlinear functions depend only on $x^1$
and only on $x^2$, respectively. Then $\mathcal C$ is the intersection of the
closures of the two groups taken separately.
\end{proposition}

\begin{proof}
After a permutation of coordinates, the graph of the block is the product
of the graphs of the two groups, and
$\conv(G^1\times G^2)=\conv G^1\times\conv G^2$: one inclusion holds
because the right side is convex, and the other because
$(\sum_i\alpha_ia_i,\sum_j\beta_jb_j)=\sum_{i,j}\alpha_i\beta_j(a_i,b_j)$.
The cone $Q$ splits in the same way, so $\mathcal C$ is the intersection of the
preimages of the two factors, and each preimage is the closure of its group
over $D$.
\end{proof}

When the rows do share variables, the gain can be strict.

\begin{example}[Two rows sharing one variable]\label{ex:two-rows}
Let $D=[-1,1]$ and consider the rows $x^2-v_1\le0$ and $-x^3+v_2\le0$.
The convex envelope of $x^2$ is $x^2$, and the convex envelope of $-x^3$ on
$[-1,1]$ takes the value $-1/4$ at $x=0$ (on $[-1/2,1]$ it is the line
through $(1,-1)$ that is tangent to $-x^3$ at $x=-1/2$). By
Corollary~\ref{cor:one-row}, the intersection of the two one-row
closures contains $(x,v_1,v_2)=(0,0,1/4)$. The aggregated cut with $a=0$
and $\lambda=(1,1)$ uses $\min_{u\in[-1,1]}(u^2-u^3)=0$ and reads
$v_1-v_2\ge0$, which this point violates. In terms of~\eqref{eq:closure},
the only probability measure on $[-1,1]$ with mean $0$ and second moment
$0$ is the point mass at $0$, whose third moment is $0$, not $1/4$.
\end{example}

That joint relaxations dominate separate envelopes is well known
\citep[Lemma~3.4]{Nowak2005}, \citep{Tawarmalani2010,LiersEtAl2021}; the
example only makes the effect concrete for the cut family used here.

\subsection{Relation to the hull of the feasible set}\label{sec:feasible-hull}

The closure $\mathcal C$ is a relaxation of the set
$\Sigma=\{v:\ Sv\in D,\ v\text{ satisfies~\eqref{eq:rows}}\}$, so
$\conv\Sigma\subseteq\mathcal C$. Equality holds when every row has its own free
direction in the remainders.

\begin{proposition}[Free remainders]\label{prop:free-remainders}
Write $v=(x,z)$ with $x$ the block, and let $M$ be the matrix formed by the
columns of $\binom AC$ that belong to $z$. If $M$ has full row rank $m+p$,
then $\mathcal C=\conv\Sigma$.
\end{proposition}

\begin{proof}
Let $v=(x,z)\in\mathcal C$. By Carath\'eodory's theorem there are $u_j\in D$ and
weights $\theta_j>0$ summing to one with $\sum_j\theta_ju_j=x$,
$\bar y:=\sum_j\theta_jg(u_j)\le b-Av$ and $\bar\eta:=\sum_j\theta_jh(u_j)=e-Cv$.
Let $M^+$ be a right inverse of $M$, let $A_x$ and $C_x$ denote the block
columns, and put
$\rho_j=\bigl(\bar y-g(u_j)+A_x(x-u_j),\ \bar\eta-h(u_j)+C_x(x-u_j)\bigr)$ and
$z_j=z+M^+\rho_j$. A direct substitution shows that $(u_j,z_j)$ satisfies
the inequality rows (because $\bar y+Av\le b$) and the equality rows, so
$(u_j,z_j)\in\Sigma$. Since $\sum_j\theta_j\rho_j=0$, the convex combination
of the points $(u_j,z_j)$ is $v$.
\end{proof}

The objective row always contains $-t$, so for the objective alone the
closure is the hull of the epigraph over $D$; this is the case studied by
\citet{LiersEtAl2021}. Without free remainders the gap can be strict.

\begin{example}[Gap to the feasible hull]\label{ex:quarter}
Let $D=[0,1]$ and take the single equality row $x^2=1/4$, whose feasible
set is $\{1/2\}$. By Corollary~\ref{cor:one-row}, $\mathcal C=\{x:x^2\le1/4\le x\}=[1/4,1/2]$,
and the two cuts $x\ge1/4$ and $x\le1/2$ suffice to describe it. The
point $x=1/4$ is the mean of the measure with mass $3/4$ at $0$ and $1/4$
at $1$, which satisfies the row in expectation.
\end{example}

This is a Lagrangian duality gap \citep{Geoffrion1974,Nowak2005}: the
rows are imposed on averages, after convexification. Even the complete
family of cuts accepts $x=1/4$, so the gap is not a failure of direction
search. It is closed by convexifying aggregated sets instead of aggregated
functions, as in the aggregation closure of
\citet{DeyMunozSerrano2022}: here $\conv\{x\in[0,1]:x^2\le1/4\}=[0,1/2]$ and
$\conv\{x\in[0,1]:x^2\ge1/4\}=[1/2,1]$ intersect in $\{1/2\}$. It is also
closed by placing the row inside $D$, which changes the support problem;
our implementation keeps nonlinear rows out of $D$ so that the support
problems stay in the classes of Section~\ref{sec:quadratic}. The hierarchy
is $\conv\Sigma\subseteq\bigcap_\lambda\conv\Sigma_\lambda\subseteq\mathcal C$, where
$\Sigma_\lambda$ is the set defined by one aggregated row, and both
inclusions can be strict. The gap can occur even when $D$ is exactly the
projection of $\Sigma$: with $D=[0,2]$, a free $z$ and the two rows
$z\le(x-2)^2/2$ and $z\le x^2/2$, the point $(x,z)=(1,1)$ lies in $\mathcal C$ (the
mean of the point masses at $0$ and $2$), whereas $z\le1/2$ on $\Sigma$.

\subsection{Exporting the eliminated row}\label{sec:export}

After the support bound is certified, the cut~\eqref{eq:agg-cut} is
assembled exactly: repeated occurrences of each original variable are
combined, giving a rational row $c\T v\ge r$. The coefficients that the
LP solver will store are binary64 numbers $\hat c$, which differ from $c$
in general.

\begin{proposition}[Safe export]\label{prop:export}
Let $\Delta=\hat c-c$, and let $L_j\le v_j\le U_j$ be valid bounds, where
$L_j$ is finite whenever $\Delta_j>0$ and $U_j$ is finite whenever
$\Delta_j<0$. Put $E=\sum_{j:\Delta_j>0}\Delta_jL_j+\sum_{j:\Delta_j<0}\Delta_jU_j$.
Then $\hat c\T v\ge\hat r$ is valid for every $\hat r\le r+E$.
\end{proposition}

\begin{proof}
$\hat c\T v=c\T v+\Delta\T v\ge r+E$.
\end{proof}

This is the bound-corrected rounding of safe LP bounds and safe cuts
\citep{NeumaierShcherbina2004,CookEtAl2009SafeCuts,EiflerGleixner2024},
applied to the eliminated row. A coefficient that is exactly representable
needs no bound, even on an unbounded variable such as the objective
variable. The correction must be applied after every transformation that
changes the row, including merging, scaling and the removal of tiny
coefficients; $\hat r$ is rounded downward and compared with $r+E$ in exact
arithmetic. If a needed bound is infinite or the correction destroys the
violation, the row is discarded; a fixed safety shift does not replace
the calculation. The bounds must be global implications of the model.
Our exporter is more conservative than Proposition~\ref{prop:export}: it
requires both bounds of a variable to be finite whenever the coefficient
of that variable changes.
